%
% Front matter
%

\begin{frontmatter}
\title{
  T2K Neutrino and Anti-Neutrino 3-Flavour Joint Analysis of  \\
  \RunPeriod (\TotalPotNuRunOneEight-POT $\nu$  \TotalPotNuBarRunOneEight-POT $\bar{\nu}$)
  data sets
}
\author[RAL,LIV]  {Costas Andreopoulos}
\author[LIV]  {Christopher Barry}
\author[WAR]  {Andrew Chappell}
\author[LANCS]  {Thomas Dealtry}
\author[RAL,LIV]  {Steve Dennis\\}
%\author[IFIC]     {Lorena Escudero S\'anchez}
%\author[WAR]{Nick Grant}
\author[GEN]     {Davide Sgalaberna}
\author[RAL,OXF]  {Raj Shah}
%\author[RAL,OXF]  {and Alfons Weber}

\address[RAL]   {Science \& Technology Facilities Council (STFC), Rutherford Appleton Laboratory,\\ Harwell Oxford Campus, OX11 0QX, UK}
\address[LIV]   {University of Liverpool, Physics Department,\\ Liverpool, L69 7ZE, UK} 
%\address[ETHZ]  {ETH Zurich, Institute for Particle Physics,\\ 8092 Zurich, Switzerland}
\address[OXF]   {University of Oxford, Subdepartment of Particle Physics,\\ Oxford, OX1 3RH, UK}
\address[WAR]   {University of Warwick, Physics Department,\\ Coventry, CV4 7AL, UK}
%\address[IFIC]  {CSIC and University of Valencia, Instituto de Fisica Corpuscular, \\ E-46071 Valencia, Spain}
\address[LANCS] {Lancaster University, Physics Department,\\ Lancaster, LA1 4YB, UK}
\address[GEN]  {University of Geneva, Departement de physique nucleaire et corpusculaire,\\ 1211 Geneve, Switzerland}

%
% Abstract
%

\begin{abstract}

We report the results of the 3-flavour, 5 sample $\nu/\bar{\nu}$ joint analysis, performed by the {\em VALOR} group, on the \RunPeriod dataset, using data taken in both neutrino and antineutrino beam mode. \color{blue}Two sets of results are reported; \cref{sec:results} presents results before and \cref{sec:results2017b} results after the inclusion of a new systematic parameter and smearing procedure to address biases associated with fake data studies involving binding energy\color{black}. The goal of this analysis is to search for CP violation by setting limits on $\delta_{CP}$ and measure the oscillation parameters $\sin^2 \theta_{13}$, $\sin^2 \theta_{23}$, $\Delta m^2_{32}$ ($\Delta m^2_{13}$). Fits are performed under both mass hierarchy hypotheses to show which is preferred by data.

The \RunPeriod dataset corresponds to an integrated J-PARC neutrino beam exposure of \TotalPotNuRunOneEight{} POT in neutrino mode and \TotalPotNuBarRunOneEight{} POT in antineutrino mode.

This analysis predicts \NskOscExpectationNumuRunOneEight{} $\pm$ \NskOscExpectationNumuErrorRunOneEight{} (syst) for neutrino $\mu$-like, \NskOscExpectationNumuBarRunOneEight{} $\pm$ \NskOscExpectationNumuBarErrorRunOneEight{} (syst) for antineutrino $\mu$-like, \NskOscExpectationNueRunOneEight{} $\pm$ \NskOscExpectationNueErrorRunOneEight{} (syst) for neutrino $e$-like, and \NskOscExpectationNueBarRunOneEight{} $\pm$ \NskOscExpectationNueBarErrorRunOneEight{} (syst) for antineutrino $e$-like single-ring \( 0 \pi \) events in Super-K and \NskOscExpectationNueCConepiRunOneEight{} $\pm$ \NskOscExpectationNueCConepiErrorRunOneEight{} (syst) \nueCConepi{} events, given the oscillation parameters of the Asimov data set A in \ref{sec:asimov_definition}. We observe respectively \NskObservedNumuRunOneEight{} for neutrino $\mu$-like, \NskObservedNumuBarRunOneEight{}  for antineutrino $\mu$-like,\NskObservedNueRunOneEight{}  for neutrino $e$-like, \NskObservedNueBarRunOneEight{}  for antineutrino $e$-like and \NskObservedNueCConepiRunOneEight{} for \nueCConepi{}-like.

\end{abstract}

\end{frontmatter}
